\scanpage{48}
% archive sheet number: 47 (upper left)
\tikz[remember picture,overlay]{\coordinate (scanXLVIIIstart) at (-.25,.25);}que la flèche\unsure{/}
\[
 H_x^i(X,\mathbb Z/n\mathbb Z)\longrightarrow
 \deleted{\unsure{H^i}}\underset{U(x)}{\underleftarrow{\text{``}\lim\text{''}}}\,
 H_c^i(U,\mathbb Z/n\mathbb Z)
\]
est un isomorphisme de pro-objets, le second membre étant représenté par le premier. On aurait alors un isomorphisme
\[
 \underline H^{-i}(Rf^!\mathbb Z/n\mathbb Z)_x
 \xrightarrow{\sim} H_x^i(X,\mathbb Z/n\mathbb Z)^\vee
\]

\tikz[remember picture,overlay]{\draw[gray,thin] (scanXLVIIIstart) to[bend right=8] (-.25,.35);}
La propriété de ``finitude'' suivante résulte aussitôt \deleted{Pro} de \deleted{ce que}\inserted{la construction} des \unsure{foncteur|foncteurs} $f_i^!$ (3.1.4.5), \inserted{et du fait qu'ils} transforment injectifs en injectifs.

\uline{Proposition} 3.1.18. (i) \uline{\deleted{\unsure{Soit|Si}}Supp\inserted{osons} le morphisme compactifiable} $f:X\longrightarrow S$ \uline{est de dimension relative} $\leq d$, \deleted{\unsure{si alors}} \deleted{\unsure{alors}}\inserted{et soit} $L\in D^+(S,\mathcal A)$, \deleted{\uline{alors}}\inserted{tel que} \deleted{si} $\underline H^i(L)=0$ \uline{pour} $i\geq k$, \uline{alors} $\underline H^i(L)=0$ \uline{pour} $i\geq k-2d$
% slanted
\margin{devrait remonter avant 3.1.5 ! ou imméd. après}

\deleted{(ii) \uline{Si} $\dim\mathrm{\,inj}(L)\leq k$, \uline{alors dim} $\mathrm{inj}(Rf^!L)\leq \unsure{kd|k1}$.}

(ii) $\dim\mathrm{\,inj}(Rf^!L)\leq\dim\mathrm{\,inj}(L)$.

% slanted
\margin{remonter également avant 3.1.6. \uline{bref}}
3.1.19. Lorsque $f$ est quasi-fini et si on fait $d=0$ dans la construction 3.1.4 du foncteur $Rf^!$, on trouve \inserted{par 3.1.3} que le foncteur $f_!$ admet un adjoint à droite $f^!$, dont $Rf^!$ est le dérivé.
% slanted
\margin{de quoi dans quoi ?}

En particulier \uline{si} $f$ \uline{est une immersion} $f:X\hookrightarrow S$, on trouve que \uline{$Rf^!$ est le dérivé du foncteur} ``\uline{sections à support dans $X$}''. En particulier
\[
 R^if^!_{\deleted{\unsure{?}}}L=\underline H_X^i(L)
 \tag*{(3.1.19.1)}
\]
\margin{\unsure{Dire} aussi que si $j$ est \unsure{éta.|étale}, $Rj^!=j^*$ ; corollaire : ``nature locale'' de $Rf^!$}

Le plan de 3.1 semble flou après 3.1.4. On pourrait prendre

a) Propriétés formelles de $Rf^!$ pour $f$ fixé : \deleted{\unsure{a)}} 3.1.5, 3.1.18, 3.1.19, 3.1.8, 3.1.9\margin{?}

b) Propriétés formelles pour $f$ variable : 3.1.6 $+$ \fbox{\shortstack[l]{chgt de base par morphisme\\ lisse (ou loc. acyclique)}}

c) Sorites sur construction locale de $Rf^!$ : 3.1.10 \unsure{à|''} 3.1.13

\diagram{\begin{tikzcd}[ampersand replacement=\&,column sep=small,row sep=small]
 X \arrow[d,dash,"f"'] \& X\times_S X \arrow[l,"p_1"'] \arrow[d,"p_2"]
 \& X \arrow[l,"\Delta"']\\
 S \& X \arrow[l,dash,"f"] \& {}
\end{tikzcd}}
\[
 \begin{aligned}
 \deleted{f}f^!&\simeq(p_1\Delta)^*f^!\simeq\Delta^*(p_1^*f^!)
 \underset{\text{chgt de base par lisse}}{\simeq}\Delta^*(p_2^!f^*)=\deleted{\unsure{\Delta^!p_2^!f^*}}\\
 &\underset{\text{pureté relative ??}}{\simeq}\Delta^!p_2^!f^*[2d](d)\simeq f^{\unsure{!|*}}[2d](d)
 \end{aligned}
\]
\margin{$\uparrow$ \uline{à ne pas} \uline{oublier} !}
% The arrow from the last note points to the boxed base-change item in b).
